% 完整中文精修稿；由冻结的 upper_zh_initial.tex 逐句修改而来。
% 需要支持中文的主文档，以及 amsmath、amssymb、amsthm 和 theorem、lemma、proposition 环境。
% 段落编号用于 upper_zh_sentence_audit.md，不排入正文。

\section{完整源线段与严格上界改进}
\label{sec:upper-constructions}

% [U001]
缩短长臂未必节省面积：其他针可能已经覆盖它失去的全部部分。
反过来，延长互补臂也未必增加面积，只要新增的整个三角形都在已有径向上边界之下。
两种效应都可利用，节省多少却必须由所得并集的面积判断。
对不受限问题最强的比较来自一个整弦操作：先延长集合内的整条弦，再以星心为中心缩回单位长度。

\begin{theorem}[严格上界构造]\label{thm:upper-chain}
存在下文定义的固定源轮廓及其实际面积代价，使得
\begin{equation}\label{eq:upper-chain}
 A_*\le c_{\mathrm{hl}}-\Gamma_{\mathrm{ms}}
 <c_{\mathrm{hl}}<c_{\mathrm{height}}<c_r<c_{\mathrm{fr}}<U,
 \qquad U=0.284167531115949,
\end{equation}
其中 $\Gamma_{\mathrm{ms}}>0$ 是固定常数。
另有一个独立构造的代价 $c_{\mathrm{new}}<c_{\mathrm{fr}}$，满足 $A_*\le c_{\mathrm{new}}$。

$c_{\mathrm{fr}},c_r,c_{\mathrm{height}},c_{\mathrm{hl}},c_{\mathrm{new}}$ 各自都是一个指定全源族的面积极限，也都能由原始有限体逼近。
对于弦归一化后的输出，只断言嵌套的上极限估计。
常数 $U$ 是已有证书确认的历史基准；本定理中后续各个差值均不在这里作数值求值。
\end{theorem}

% [U003]
式~\eqref{eq:upper-chain} 中的符号表示不同的集合或不同的极限，并非同一个数的逐次小数近似。
特别地，论证不使用 $c_{\mathrm{new}}$ 与 $c_r$ 或后续代价之间的比较。
我们先构造实际面积极限，再比较互补臂的代价，最后使正面积节省在两次离散化后仍然保留。

\subsection{周期记录包含整条线段}
\label{subsec:upper-periodic}

% [U004]
设 $p:[0,1]\to(0,1)$ 连续且满足 $p(1-t)=1-p(t)$；设 $h:[0,1]\to[0,\infty)$ 连续，满足 $h(1-t)=h(t)$ 及 $h(0)=h(1)=0$。
对 $y=j+t$、$j\in\mathbb Z$、$0\le t\le1$，令
\begin{equation}\label{eq:upper-periodic-data}
 P(y)=\frac12+(-1)^j\left(p(t)-\frac12\right),
 \quad M(y)=1-P(y),\quad A(y)=(-1)^{j+1}h(t).
\end{equation}
端点恒等式保证这些定义在整数接缝处一致。
这些函数以二为周期，并满足 $P(y+1)=M(y)$、$A(y+1)=-A(y)$。
对奇数 $n$，令 $\delta_n=\pi/n$，并定义
\begin{equation}\label{eq:upper-actual-source}
 N_{n,y}=\{xu_{\delta_n y}+\delta_n A(y)n_{\delta_n y}:
                 -M(y)\le x\le P(y)\},\qquad
 E_n[p,h]=\bigcup_{0\le y\le n}\operatorname{conv}(0,N_{n,y}).
\end{equation}
由 $P+M=1$，每条底边的长度都是一，而其方向覆盖 $[0,\pi]$。
交替放置并不要求加入对径集合 $-E_n$。

% [U005]
沿底边点向原点填充，就在同一角度上得到所有更小的半径。
这里的相位是将实际角度按尺度 $\delta_n$ 展开后取模二的坐标；有限角度的精确公式见下面的证明。
固定极限半径 $\rho>0$ 后，要得到完整截面，就须让纵向坐标从 $\rho$ 遍历到臂端点，而非只取端点值。
在带有总质量为二的长度测度的 $\mathbb T_2=\mathbb R/2\mathbb Z$ 上，定义
\begin{align}
 S_+^{p,h}(\rho)&=\bigcup_{\substack{0\le y\le2\\P(y)\ge\rho}}
 \{y+A(y)/x\pmod2:\rho\le x\le P(y)\},\\
 S_-^{p,h}(\rho)&=\bigcup_{\substack{0\le y\le2\\M(y)\ge\rho}}
 \{y-A(y)/x\pmod2:\rho\le x\le M(y)\}.
\end{align}
这个完整源族的\emph{实际极限代价}为
\begin{equation}\label{eq:upper-physical-cost}
 C_{\mathrm{phys}}(p,h)=\frac\pi2\int_0^\infty
 \rho\bigl(|S_+^{p,h}(\rho)|+|S_-^{p,h}(\rho)|\bigr)\,d\rho.
\end{equation}
这是并集的积分，不是源三角形面积之和，也不是端点包络的替代量。
关系 $S_-+1=S_+$ 说明二者的相位长度相同，并不把实际平面中的两个半平面等同起来。

\begin{proposition}[两阶段实际面积恢复]\label{prop:upper-recovery}
对本节使用的每一对固定轮廓，即冻结的多项式对及下文构造的有限分片代数或多项式修改，均有
\[
 \lim_{\substack{n\to\infty\\n\ \mathrm{odd}}}|E_n[p,h]|
       =C_{\mathrm{phys}}(p,h).
\]
对每个固定奇数 $n$，存在原始有限体 $K_{n,k}[p,h]$，满足
\[
 \lim_{k\to\infty}|K_{n,k}[p,h]|=|E_n[p,h]|.
\]
网格尺度趋于零，并细分所有源分片边界；每个闭小源区间上，都在同一个代表参数处取 $b=(-1)^j(p-1/2)$ 和带符号高度 $\delta_n(-1)^{j+1}h$ 的值，再同时冻结。
并集中保留接缝两侧的所有闭记录。
\end{proposition}

\begin{proof}
% [U006]
令 $H=\max_{[0,2]}|A|$，并固定满足 $\delta_nH<\rho$ 的 $\rho>0$。
正臂的精确定义域与相位为
\[
 \sqrt{\rho^2-\delta_n^2A(y)^2}\le x\le P(y),\qquad
 y+\delta_n^{-1}\arctan\bigl(\delta_n A(y)/x\bigr).
\]
负臂使用 $M$，并将反正切前的符号反转。
因此，在取极限之前，真实半径和底边的完整内部都已保留。
避开 $P$ 和 $M$ 的有限多个临界值、常值分片值及接缝值后，所有激活端点都是横截的。
这些连通定义域在上述扰动下持续存在，且对充分大的 $n$，其上有 $x\ge\rho/2$。
相位一致收敛到 $y\pm A(y)/x$。
连通定义域的实相位像是完整区间；极值收敛先给出各区间长度收敛，再给出它们模二后的有限并集长度收敛。

% [U007]
在这个固定正半径上，相位位移有界，所以把无限周期源截为 $[0,n]$ 只会改变有界个数的首尾单元。
这些首尾单元的实际角宽趋于零。
分别计入正负两侧的实际区域，得到极限角长度 $\frac\pi2(|S_+|+|S_-|)$。
所有集合都在同一个圆盘中，每个圆周截面的角长度至多为 $2\pi$。
半径为 $\rho_0$ 的圆盘把整个小半径部分的贡献控制在 $\pi\rho_0^2$ 以内，再由控制收敛定理得到式~\eqref{eq:upper-physical-cost}。

% [U008]
固定 $n$ 后，端点的一致逼近使 $K_{n,k}$ 落在紧集 $E_n$ 逐渐缩小的外邻域内，从而得到面积的上极限估计。
对于下极限，每个非退化源三角形的内点都在其冻结近似中持续存在。
逐个考察闭解析源坐标片，也包括代数端点处的收敛 Puiseux 坐标片：端点曲线 $e$ 要么只有有限多个 $\det(e,e')$ 的零点，要么位于一条固定射线上。
在非临界的侧边点处，略微改变源参数，就能把该点的严格径向缩点移入相邻三角形的内部。
例外的侧边射线只有有限条；在其余每条射线上，只有最外端点可能不属于三角形内部的并集。
有限条例外射线的面积为零；在其余每条射线上，遗漏的至多一个最外端点也有一维径向测度零。极坐标 Fubini 定理因而说明整个遗漏集的面积为零。
记非退化源三角形内部的并为 $\mathcal I_n$，则 $|\mathcal I_n|=|E_n|$，且 $\mathbf1_{\mathcal I_n}\le\liminf_k\mathbf1_{K_{n,k}}$ 逐点成立；Fatou 引理给出面积的下极限估计。
正高度源端点始终保留其完整三角形；测度论证中略去的只有上述边界射线。
由此得到的是实际面积收敛；论证没有借 Hausdorff 收敛代替面积收敛。
\end{proof}

% [U009]
冻结的轮廓对为 $p_0=1/2+\alpha_{\mathrm{fr}}$ 与 $s$；本节末尾列出 $\alpha_{\mathrm{fr}}$ 和 $s$ 按幂次递增排列的精确系数。
定义 $c_{\mathrm{fr}}=C_{\mathrm{phys}}(p_0,s)$。
电子符号证书与包络证书给出了下文使用的严格凹性、反射关系、开盒约束、两个严格起始不等式及递增端点映射，也给出了历史比较 $c_{\mathrm{fr}}<U$。
附录~\ref{sec:formal-evidence} 按数学输入列出这些证据的对应关系；它们只支持上述固定输入，不能推广为所有同次数多项式都满足这些不等式。

\subsection{互补臂产生凸铰链函数}
\label{subsec:upper-hinge}

% [U010]
固定 $s\in C^2([0,1])$，满足 $s(1-t)=s(t)$、$s(0)=s(1)=0$，且在 $[0,1]$ 上有 $s''<0$。
满足反射关系的放置函数 $p\in C^1([0,1])$ 称为\emph{起始可容许}，若它在整个闭区间上满足 $0<p<1$，并且满足
\begin{equation}\label{eq:upper-birth}
 p(t)>\frac{s(t)}t,\qquad
 1-p(t)>\frac{s(t)}{1-t}\qquad(0<t<1).
\end{equation}
令 $d=s'(0)$、$g(t)=t-s(t)/p(t)$。第一个起始不等式给出 $g(t)>0$，第二个对反射臂给出同样的结论：两条臂的端点相位都已经越过各自的零相位。再令
\begin{equation}\label{eq:upper-tail-threshold}
 T_s(z)=\max_{0\le u\le1}\frac{s(u)}{u+z}\quad(z>0),
 \qquad B_s(t)=\max_{0\le u\le1}\frac{s(t)+s(u)}{t+u}\quad(t>0),
\end{equation}
其中在零点以连续延拓定义 $T_s(0)=B_s(0)=d$。
$T_s(z)$ 描述相位 $-z$ 上的径向上边界：在内部源 $0<u<1$ 处解 $u-s(u)/x=-z$，得到 $x=s(u)/(u+z)$。起始条件保证这个底边坐标合法。
因此，只要高度轮廓相同，每个起始可容许放置都包含同一个负相位尾部；反射后，它也参与正相位上的比较。

\begin{lemma}[阈值与完整径向上边界]\label{lem:upper-roof}
对 $0<t\le1$，定义 $B_s(t)$ 的最大值由唯一的源参数 $u_t\in(0,\min(t,1/2))$ 取得，且
\[
 B_s(t)=s'(u_t),\qquad
 B_s'(t)=\frac{s'(t)-B_s(t)}{t+u_t},\qquad
 d>B_s(t)>\frac{s(t)}t>s'(t).
\]
对任意起始可容许的 $p$，定义
\[
 \tau(z)=\min\{t\in[0,1]:g(t)\ge z\},\qquad 0<z<1.
\]
完整正向径向上边界为 $R(z)=p(\tau(z))$，实际代价为
\begin{equation}\label{eq:upper-roof-cost}
 C_{\mathrm{phys}}(p,s)=\pi\int_0^1\max\{R(z),T_s(z)\}^2\,dz.
\end{equation}
若臂长 $x>s(t)/t$，则它与互补尾部的比较满足
\[
 x>T_s(t-s(t)/x)\quad\Longleftrightarrow\quad x>B_s(t).
\]
将两边都改成等号或都反转不等号，同样等价。
\end{lemma}

\begin{proof}
% [U011]
阈值商对源参数求导后的分子为 $(t+u)s'(u)-s(t)-s(u)$。
这个分子对 $u$ 的导数为 $(t+u)s''(u)<0$，在 $u=0$ 和 $u=t$ 处符号相反。
它在 $1/2$ 处为负，因而极大源具有所述位置且唯一；随后隐式求导给出各恒等式。
与 $T_s$ 的比较来自恒等式
\[
 x(u+t-s(t)/x)-s(u)=x(t+u)-s(t)-s(u),
\]
这里使用了分母为正以及最大值能够取到。

% [U012]
起始条件给出 $g(0)=0$、$g(1)=1$，并在内部给出 $0<g(t)<t$，所以 $\tau(z)$ 存在且 $g(\tau(z))=z$。
一条完整正底边只有在满足 $g(u)\ge z$ 的源参数 $u$ 处才能到达相位 $z$，对应半径为 $s(u)/(u-z)$。
在第一个这样的源参数处，导数分子 $(u-z)s'(u)-s(u)$ 为负，此后还会严格递减。
所以，即使合法源集合不连通，每个更晚的合法源给出的半径也都更小。
这就在不丢弃底边内点的前提下证明了首次到达的径向上边界公式。
负相位给出 $T_s$，反射给出另一侧的提升；由于 $T_s$ 递减，第一个单元中每个更远的等价相位表示也都受 $T_s(z)$ 控制。
第二个单元由相应反射得到，再计入实际平面的两侧，就得到式~\eqref{eq:upper-roof-cost}。
\end{proof}

% [U013]
定义铰链函数及其反射配对形式为
\begin{align}
 \Phi_s(t,x)&=(x-B_s(t))_+^2
             +2(B_s(t)-s'(t))(x-B_s(t))_+,\\
 \Psi_t(x)&=\Phi_s(t,x)+\Phi_s(1-t,1-x),
 \qquad m_s(t)=\min_{0\le x\le1}\Psi_t(x),\\
 L(s)&=\pi\left(\int_0^1T_s(z)^2\,dz+
                         \int_0^{1/2}m_s(t)\,dt\right).
\end{align}
这里 $v_+=\max(v,0)$。
铰链函数是凸的，但一般在 $x=B_s(t)$ 处不可微。
线性项不可省略：零代价臂一旦延伸越过阈值，就立即增加暴露面积。

\begin{theorem}[铰链比较及几何等号条件]\label{thm:upper-hinge}
设 $p$ 起始可容许，且 $\min_{[0,1]}g'>0$。
则
\begin{equation}\label{eq:upper-hinge-identity}
 C_{\mathrm{phys}}(p,s)=\pi\left(\int_0^1T_s^2+
                                      \int_0^1\Phi_s(t,p(t))\,dt\right)
 \ge L(s).
\end{equation}
对 $0\le t\le1/2$，令 $b=B_s(t)$、$\bar b=B_s(1-t)$、$\sigma=s'(t)$。
若 $b+\bar b\ge1$，则 $\Psi_t$ 的极小点集为 $[1-\bar b,b]$，最小值为零。
若 $b+\bar b<1$，则其唯一极小点为
\[
 x_t=\operatorname{clamp}\left(\frac12+\sigma,[b,1-\bar b]\right),
\]
其中对 $v\le w$，定义 $\operatorname{clamp}(x,[v,w])=\min(w,\max(v,x))$。
对于已经满足本定理几何假设的放置，等号成立当且仅当每个 $t\in[0,1/2]$ 处的 $p(t)$ 都属于相应极小点集。
\end{theorem}

\begin{proof}
% [U014]
令 $H(z)=\int_0^zT_s(w)^2\,dw$，并在 $x>B_s(t)$ 时定义
\[
 F(t,x)=s(t)(x-B_s(t))-H(t-s(t)/x)
                         +H(t-s(t)/B_s(t));
\]
在 $x\le B_s(t)$ 时把 $F$ 延拓为零。
由阈值恒等式，拼接曲线上 $F$ 及其两个一阶偏导数都为零。
因此，链式法则穿过任意接触集仍然成立，并给出
\begin{equation}\label{eq:upper-primitive}
 (p^2-T_s(g)^2)_+g'
       =\frac{d}{dt}F(t,p(t))+\Phi_s(t,p(t)).
\end{equation}
沿曲线的导数有界，且 $F$ 在两端都趋于零，故其导数积分不产生边界项或奇异项。
在式~\eqref{eq:upper-roof-cost} 中通过递增映射 $g$ 换元，即得恒等式。

% [U015]
第一个激活铰链的导数是 $2(x-\sigma)$，第二个的导数是 $-2(1-x+\sigma)$。
如果阈值区间重叠，两个铰链都在共同区间上为零。
两阈值之间若留有间隙，配对导数在间隙内为 $4x-2-4\sigma$；在间隙外向间隙移动，函数值总会下降。
由此得到上述截断极小点，也包括边界情形。
最后，反射关系将完整铰链积分写成 $\int_0^{1/2}\Psi_t(p(t))\,dt$，逐点比较便证明了下界。
这里没有把几何可行函数上的下确界移入积分。
\end{proof}

% [U016]
$L(s)$ 给出这个源类的松弛下界，是否能作为最小值取到还要检验几何条件。
拟议的极小放置仍须按要求的正则性拼接，满足起始条件，并满足
\[
 J_p(t):=p(t)^2-p(t)s'(t)+s(t)p'(t)>0,
 \qquad g'=J_p/p^2.
\]
逐点凸极小化不会自动保证这个微分条件。

\begin{proposition}[松弛最小值不可达的非空类]
\label{prop:upper-unattainable}
对每个 $\kappa\in(0,1/4)$，选择 $\varepsilon>0$ 使得
\[
 \kappa^2\frac{\sqrt{\varepsilon^2+1/4}-\varepsilon}{\varepsilon}>\frac14,
\]
并令
\[
 s(t)=\kappa\left(\sqrt{\varepsilon^2+1/4}
                -\sqrt{\varepsilon^2+(t-1/2)^2}\right).
\]
定理~\ref{thm:upper-hinge} 中的类非空，但其中 $C_{\mathrm{phys}}(p,s)$ 的下确界严格大于 $L(s)$。
\end{proposition}

\begin{proof}
% [U017]
这里 $|s'|<\kappa$ 且 $B_s\le d<\kappa$，所以 $p=1/2$ 满足起始条件以及 $g'>1-2\kappa>0$。
在整个区间上，唯一的逐点极小放置为 $p_*=1/2+s'$。
它满足起始条件和反射关系，但
\[
 J_{p_*}(1/2)=\frac14-
 \kappa^2\frac{\sqrt{\varepsilon^2+1/4}-\varepsilon}{\varepsilon}<0.
\]
此外，每个配对铰链在盒区间上的分布二阶导数至少为二，故代价趋近 $L(s)$ 会迫使 $p_j\to p_*$ 于 $L^2$。
取一个几乎处处收敛的子列后，非递减函数 $g_j=t-s/p_j$ 也会几乎处处收敛到 $g_*=t-s/p_*$。
这种极限有一个非递减代表元，而连续函数 $g_*$ 在 $1/2$ 附近严格递减。
这一矛盾排除了固定高度、端点映射递增的类中的任何恢复序列。
\end{proof}

\subsection{折叠对应精确的实际面积余项}
\label{subsec:upper-fold}

\begin{theorem}[有序交叉余项]\label{thm:upper-fold}
沿用第~\ref{subsec:upper-hinge} 节中 $s\in C^2([0,1])$ 的严格凹性、反射关系及零端点条件，并设 $p\in C^1([0,1])$ 起始可容许。这里仅撤去 $\min_{[0,1]}g'>0$ 的要求，允许 $g'$ 改变符号。
对几乎每个 $z\in(0,1)$，它的根依次为
\[
 t_1<\cdots<t_{2k+1},\qquad g(t_i)=z,
\]
导数符号按 $+,-,+,\ldots,+$ 交替。
定义 $w_i(z)=\max\{p(t_i)^2,T_s(z)^2\}$，并定义
\begin{equation}\label{eq:upper-fold-def}
 D_{\mathrm{fold}}(p,s)=\pi\int_0^1
               \sum_{j=1}^{k}\bigl(w_{2j}(z)-w_{2j+1}(z)\bigr)\,dz.
\end{equation}
则
\begin{equation}\label{eq:upper-fold-identity}
 C_{\mathrm{phys}}(p,s)-L(s)=
 \pi\int_0^{1/2}\bigl(\Psi_t(p(t))-m_s(t)\bigr)\,dt
                 +D_{\mathrm{fold}}(p,s),\qquad D_{\mathrm{fold}}\ge0.
\end{equation}
此外，$D_{\mathrm{fold}}=0$ 当且仅当在开激活集 $\{t\in(0,1):p(t)>B_s(t)\}$ 上有 $g'\ge0$。
这一公式也适用于连续 Lipschitz 放置，只要它由有限个闭 $C^1$ 分片组成，并在各分片的激活内部使用同一导数条件。
\end{theorem}

\begin{proof}
% [U018]
临界集 $\{g'=0\}$ 的像的长度为零：用 $\{|g'|<\eta\}$ 覆盖它，在这个开集的区间分支上估计像的长度，再令 $\eta$ 递减到零。
非临界值处的根孤立，由紧性知根只有有限个；又由 $g(0)<z<g(1)$，其导数符号交替。
引理~\ref{lem:upper-roof} 还给出更强的几何次序
\[
 p(t_1)>p(t_2)>\cdots>p(t_{2k+1}).
\]
将首次到达半径与尾部半径取大后平方，再减去带符号交叉和，恰为
\[
 w_1-\sum_{i=1}^{2k+1}(-1)^{i+1}w_i
              =\sum_{j=1}^k(w_{2j}-w_{2j+1})\ge0.
\]
保证余项非负的是这个半径次序，无须给重叠重数设上界。

% [U019]
令 $M_s=\max(\|p\|_\infty,d)$；逐分支换元将带绝对值的加权交叉积分控制在 $M_s^2\int_0^1|g'|$ 以内。
因此，即使折叠次数没有统一上界，带符号计算仍然可积。
恒等式~\eqref{eq:upper-primitive} 是逐点链式法则，不要求 $g'>0$。
此时尾部项通过 $H(g(1))-H(g(0))$ 积分，而不是通过保向换元积分。
这些恒等式证明了式~\eqref{eq:upper-fold-identity}。
各光滑分片的原函数在接缝处取值相同，内部边界项因而逐一抵消；平坦分片和接缝处都不产生原子项。

% [U020]
如果每个向下交叉都未激活，它的半径及其后向上交叉处更小的半径都不超过 $T_s$，这一对的贡献便为零。
反之，若某个激活点处有 $g'<0$，连续性就给出一个激活的递减区间，其像具有正长度。
该像中几乎每个相位都有一个严格高于尾部的向下交叉，其后跟着半径更小的向上交叉，因此对应差值为正。
这就证明了零余项判据的两个方向。
命题~\ref{prop:upper-unattainable} 中的针族在这个扩大类中是合法的，但中央的激活折叠使 $D_{\mathrm{fold}}(p_*,s)>0$。
单个放置的余项为正，并不能推出 $L(s)$ 与扩大类的下确界之间存在正间隙。
\end{proof}

% [U021]
从同一个参考放置 $p_0$ 出发，可以作两种朝向铰链极小点的投影；它们给出的放置不同。
记 $B=B_s$，定义
\begin{align}
 r(t)&=\operatorname{median}\bigl(B(t),p_0(t),1-B(1-t)\bigr),\label{eq:upper-band}\\
 q(t)&=\begin{cases}
 \operatorname{clamp}(1/2+s'(t),[B(t),1-B(1-t)]),&B(t)+B(1-t)<1,\\
 r(t),&B(t)+B(1-t)\ge1.
 \end{cases}\label{eq:upper-completed}
\end{align}
这里的中位数是三个实数参数中居中的一个，不预设两个阈值的大小次序。
$q$ 和 $r$ 都满足反射关系及开盒约束，都是 Lipschitz 有限分片放置，并满足严格内部起始条件。
事实上，若 $\ell=s(t)/t$、$V=1-s(t)/(1-t)$，则 $\ell<p_0<V$、$B>\ell$、$1-B(1-t)<V$；无论两个阈值怎样排序，投影到它们之间都保留这两个严格不等式。
两阈值相遇时，投影区间收缩为单点，这也保证了交界处的连续性。

% [U022]
补全后的选择函数满足
\[
 C_{\mathrm{phys}}(q,s)=L(s)+D_{\mathrm{fold}}(q,s),
\]
命题~\ref{prop:upper-recovery} 又将这个实际代价恢复为有限体的面积极限。
它的激活折叠余项是否为零，仍是开放问题。
未补全选择函数 $\operatorname{median}(B,1/2+s',1-B(1-t))$ 的内部起始条件仍未解决；已经证明的 $q$ 的起始性质不能移用到它身上。
这两个未决问题都不妨碍对另一个放置 $r$ 作严格的实际面积比较。

\subsection{冻结的高次放置为何不是驻点}
\label{subsec:upper-polynomial}

\begin{theorem}[多项式接触刚性与下降]\label{thm:upper-polynomial}
设多项式 $s$ 和 $p$ 满足定理~\ref{thm:upper-hinge} 的全部假设，并设 $\deg p>\deg s'\ge1$。
则在满足反射关系、起始可容许且 $g$ 的导数有正下界的 $C^1$ 类中，$p$ 不是 $C_{\mathrm{phys}}(\cdot,s)$ 的局部极小点。
存在任意小的 $C^1$ 扰动，具有紧内部支集和有限个多项式分片，并严格降低实际代价。
对于冻结轮廓对，$\deg p_0=19$、$\deg s'=15$，所以固定其中一个扰动即可定义
\[
 c_{\mathrm{new}}:=C_{\mathrm{phys}}(p_0+\epsilon v,s)<c_{\mathrm{fr}}.
\]
\end{theorem}

\begin{proof}
% [U023]
假设次数大于 $\deg s'$ 的多项式 $Q$ 在某个开区间上与 $B_s$ 一致。
阈值恒等式在该区间给出
\[
 \mathcal R(t)=\frac{s'(t)-Q(t)}{Q'(t)}-t=u_t,
 \qquad s'(\mathcal R(t))=Q(t).
\]
清除分母后，这成为复数域上的有理恒等式。
若约分后的 $\mathcal R$ 有有限极点，将它代入非常数多项式 $s'$，就会产生一个无法抵消的最高阶极点。
所以 $\mathcal R$ 是多项式。
若 $m=\deg Q$，它在无穷远处的首项行为为 $-(1+1/m)t$，因此次数为一。
它与 $s'$ 的复合于是具有次数 $\deg s'$，与 $\deg Q>\deg s'$ 矛盾。
反射同样排除了它与 $1-B_s(1-t)$ 在开区间上接触的可能。

% [U024]
由 $g'(0)>0$ 得 $p(0)>d=B_s(0)$，故零点附近有一个第一个铰链激活的区间。
去掉第二个铰链的闭接触集以及 $p-1/2-s'$ 的有限多个零点，选取紧区间 $I=[a_0,b_0]\subset(0,1/2)$，使铰链激活状态固定，且
\[
 D(t):=\partial_x\Psi_t(p(t))
 =\begin{cases}2(p-s'),&\text{仅第一个铰链激活},\\
                 4p-2-4s',&\text{两个铰链都激活}
   \end{cases}
\]
与零保持正距离。
在 $I$ 上令 $w(t)=(t-a_0)^2(b_0-t)^2$，在 $I$ 外令其为零。
在 $I$ 上取 $v=-\operatorname{sign}(D)w$，按 $v(1-t)=-v(t)$ 反射延拓，其余处置零。
端点的二重零点使 $v$ 在全局为 $C^1$。
扰动方向必须按实际配对导数选取，不能一律按 $p-1/2-s'$ 选取。

% [U025]
对所有充分小的正 $\epsilon$，起始条件、开盒约束、激活状态和导数条件的严格裕量同时保留。
在这个未改变的激活区域内，精确的有限差分为
\[
 C_{\mathrm{phys}}(p+\epsilon v,s)-C_{\mathrm{phys}}(p,s)
     =\pi(\epsilon L_0+\epsilon^2Q_0),
\quad L_0=-\int_I|D|w<0,\quad Q_0=\int_I\nu w^2>0,
\]
其中单激活时 $\nu=1$，双激活时 $\nu=2$。
固定任意满足 $0<\epsilon<-L_0/Q_0$ 的可容许 $\epsilon$。
这样在 $n$ 和记录网格变化之前就固定了一个正的面积节省，即旧成本减去新成本，命题~\ref{prop:upper-recovery} 再把它传递到原始有限体。
\end{proof}

\subsection{带投影、被遮住的高度增量与第二次纵向下降}
\label{subsec:upper-profile-chain}

\begin{proposition}[固定带投影严格节省面积]
\label{prop:upper-band}
对于精确冻结的 $p_0,s$，式~\eqref{eq:upper-band} 中的放置 $r$ 满足
\[
 c_r:=C_{\mathrm{phys}}(r,s)<c_{\mathrm{fr}}.
\]
\end{proposition}

\begin{proof}
% [U026]
在源参数 $t$ 处，简记 $b=B_s(t)$、$v=1-B_s(1-t)$。
若 $b<v$，从下方投影只把正臂延长到 $b$，从上方投影只把负臂延长到 $1-v=B_s(1-t)$。
若 $v\le b$，从下方投影只把正臂延长到 $v\le b$，从上方投影只把负臂延长到 $1-b\le B_s(1-t)$。
四种情形都使每条被延长的臂止于自身阈值或其下方。
考察被延长的正臂中 $0<X\le B_s(t)$ 的坐标，令 $z=t-s(t)/X$。若 $t>0$ 且 $z\le0$，则 $X\le s(t)/t<p_0(t)$，这个底边点本来就在旧线段内；$t=0$ 时，$X\le B_s(0)=d\le p_0(0)$ 也给出同一包含。若 $z>0$，阈值等价关系给出 $X\le T_s(z)$。这些点均已属于旧的实际并集。
特别地，$0<t<1$ 时的相位 $z=0$ 给出 $X=s(t)/t>0$，并不表示纵向坐标 $X=0$；这样的底边点已在旧线段内，径向填充也仍在旧三角形内。
尾部已经作径向填充，因此新增的整个三角形都被遮住，而不仅是端点。
合并实际的反射提升及等价相位表示后，就得到全局实际面积不增。

% [U027]
附录~\ref{sec:formal-evidence} 所列的固定带区间支配输入，在整个源区间上给出
\[
 I_0=(2^{-20},2^{-19}),\qquad
 p_0(t)\ge B_s(t),\quad p_0(t)\ge1-B_s(1-t).
\]
证书在这里提供的是整个区间上的断言：端点映射将 $I_0$ 及其反射送入已认证的首段、末段支配范围，那里的驻定尾部候选又恰好等于 $T_s$。
多项式接触刚性使这两个不等式在某个非空开子区间 $J_0\subset I_0$ 上同时严格成立。
在那里 $r(t)<p_0(t)$，且原来的首次到达端点严格高于尾部。
由引理~\ref{lem:upper-roof}，每个更晚的合法底边坐标都严格更低；每条新增臂又都在共同尾部内。
因此，新集合失去了这个暴露端点。

% [U028]
一个端点缺失还不等于面积下降。先在低于该端点的某个正半径处截断相位--半径模型。
包括径向参数和有限多个相关等价相位表示在内，完整源定义域都是紧的，所以新并集在此处是闭的。
这个缺失点的某个邻域与新集合不交，旧的连续径向上边界却从下方填满其中一个非空开集。
其面积密度为 $\rho\,d\rho\,dz$ 的正倍数，因而失去的面积为正。
结合全局不增，就证明了严格比较。
\end{proof}

% [U029]
对于这个实际的 $r$，引理~\ref{lem:upper-roof} 的导数公式说明 $B_s(t)$ 与 $1-B_s(1-t)$ 都严格递减。冻结多项式 $p_0$ 的非递增性来自附录~\ref{sec:formal-evidence} 所列的单调性输入；它又是非常数多项式，若在两个不同参数处取值相同，非递增性就迫使它在两者之间恒定，进而成为常数多项式，矛盾。因此 $p_0$ 也严格递减。
对任意两个递增排列的参数，三个函数各自的下降量都有一个共同的正下界。中位数对每个坐标单调，且三坐标同时减去同一个数时中位数也减去该数，所以 $r$ 严格递减。
令 $P=r(0)$，仍记 $d=s'(0)$。
同一附录所列的端点与尾部界给出
\[
 d=\frac{107161751}{128520000}<\frac{21}{25},
 \qquad T_s(1)<\frac4{25},\qquad P>d.
\]
最后一个不等式来自 $B_s(0)=d$、$B_s(1)=T_s(1)$ 及 $p_0(0)>d$。
该尾部包围估计用的是输出相位 $1/2$ 对应的实际驻定源区间，以及冻结高度 $s(1/2)$；源参数中点不能直接当作这个输出相位。

\begin{proposition}[被端点上边界遮住的高度增量]
\label{prop:upper-height}
存在 $0<\alpha<\beta<1/2$ 及固定的 $\epsilon_h>0$，使得令
\[
 \psi(t)=\begin{cases}(t-\alpha)^3(\beta-t)^3,&\alpha\le t\le\beta,\\0,&\text{其余情形},\end{cases}
 \quad h_1(t)=\psi(t)+\psi(1-t),\quad a=s+\epsilon_h h_1,
\]
后，轮廓对 $(r,a)$ 保持反射关系、严格凹性和起始条件，且
\[
 c_{\mathrm{height}}:=C_{\mathrm{phys}}(r,a)<c_r.
\]
区间与振幅都预先固定，不依赖之后的奇数 $n$ 和有限网格。
\end{proposition}

\begin{proof}
% [U030]
对 $0<\rho<P$，$r$ 的严格递减性使激活源集成为一个前缀 $[0,\tau_\rho]$。
若 $\rho\le r(1)$，则 $\tau_\rho=1$；否则 $r(\tau_\rho)=\rho$。
以下 $b$ 取 $s$ 或 $a$；两者均对称、严格凹，端点为零，并满足起始条件。定义
\[
 L_b(\rho)=\max_{0\le t\le\tau_\rho}\bigl(b(t)/\rho-t\bigr),
 \qquad G_b(\rho)=\max_{0\le t\le\tau_\rho}\bigl(t-b(t)/r(t)\bigr).
\]
完整相位定义域的底部是连通前缀，每条竖直纤维也连通，因而相位像是整个区间 $[-L_b,G_b]$。
令 $M_b(\rho)=\max(L_b(\rho),G_b(\rho))$。
反射并集模二后的长度为 $2\min(1,M_b(\rho))$，因此
\begin{equation}\label{eq:upper-height-cost}
 C_{\mathrm{phys}}(r,b)=2\pi\int_0^P
      \rho\min(1,\max(L_b(\rho),G_b(\rho)))\,d\rho.
\end{equation}
定义 $L_b$ 的最大值也可在整个 $[0,1]$ 上取。内部极大源满足 $b'(t)=\rho$；严格凹性与 $b(0)=0$ 给出 $b(t)/t>b'(t)$，起始条件再给出 $r(t)>b(t)/t$，所以它已被激活。边界极大值只能在 $t=0$ 取得，而 $\rho<P=r(0)$ 保证它也被激活。

% [U031]
由 $P>d$，端点映射 $g_{r,s}=t-s/r$ 在足够短的起始邻域上，其几乎处处导数有一个正下界。
在 $\rho=d$ 处尾部最大值为零，而该邻域内一个足够小的正源参数给出 $G_s(d)>0$。
选闭窗口 $\mathcal N=[d-\eta,d+\eta]\subset(0,P)$，使其上有 $G_s-L_s>0$，并定义
\[
 \kappa_h=\frac14\min_{\rho\in\mathcal N}(G_s(\rho)-L_s(\rho))>0.
\]
在更小的起始邻域内选 $[\alpha,\beta]$，使 $r>d+\eta$、$s'$ 严格落在 $\mathcal N$ 内，且端点映射递增。
把 $\epsilon_h$ 选得足够小，使凹性、起始条件及该邻域上的递增性保留，使支集上的 $a'$ 留在 $\mathcal N$ 内，并满足
\[
 \epsilon_h\|h_1\|_\infty
       \left(\frac1{\min r}+\frac1{d-\eta}\right)<2\kappa_h.
\]
有限个紧集上的正裕量条件可同时满足，允许振幅因此包含一个非空区间。

% [U032]
对每个半径，由 $a\ge s$ 得 $G_a\le G_s$。
若 $L_a(\rho)>L_s(\rho)$，新的极大源只能落在正凸起内部；它满足 $a'(t)=\rho$，所以 $\rho$ 必在 $\mathcal N$ 内。
反射支集位于 $t>1/2$。总高度 $a$ 对称且严格凹，故其在这里的导数 $a'(t)<0$，不能满足正半径所需的极大值方程 $a'(t)=\rho>0$。这里不要求附加凸起自身的导数处处为负。
在 $\mathcal N$ 上，上述振幅约束保持 $G_a>L_a$，所以变化尾部的一切增量都被遮住。
在 $\mathcal N$ 外，尾部不变，端点占据也不增加。
这样，完整实际截面在所有半径上都不增加，极大源改变也不例外。

% [U033]
选一个非退化闭区间 $J\subset(\alpha,\beta)$。
在具有正宽度的半径区间 $r(J)$ 上，激活端点属于 $J$，两个端点映射都在整个激活前缀上递增，两个尾部都为零，且截面尚未饱和。
其损失恰为 $\epsilon_h\psi(\tau_\rho)/\rho$，故
\[
 c_r-c_{\mathrm{height}}
 \ge2\pi\epsilon_h\int_{r(J)}\psi(\tau_\rho)\,d\rho>0.
\]
区间的正宽度只需 $r$ 严格递减，不需要另给 $-r'$ 一个正下界。
\end{proof}

\begin{proposition}[固定新高度后的纵向下降]
\label{prop:upper-longitudinal}
对于刚构造并固定的 $a$，存在反射反对称的 $C^1$ 有限分片函数 $v_h$，其支集严格位于已有高度支集区间内部，并存在固定 $\eta_h>0$，使得
\[
 p_*=r+\eta_hv_h,\qquad
 c_{\mathrm{hl}}:=C_{\mathrm{phys}}(p_*,a)<c_{\mathrm{height}}.
\]
放置 $p_*$ 仍然严格递减、满足开盒约束和起始条件，并保留原端点值。
\end{proposition}

\begin{proof}
% [U034]
在已有高度支集的左边界 $\alpha$ 附近，有 $r>d>B_a$。
因此，有限分片中位数在 $\alpha$ 的某个右邻域上选取 $1-B_s(1-t)$ 或 $p_0$。
这就分成以下两种情形；纵向支集都在已经固定的高度凸起内部选取。
正的高度增量在其整个内部给出 $1-B_a(1-t)<1-B_s(1-t)$。

% [U035]
若 $r=1-B_s(1-t)$，在这个分支内选紧区间 $I=[a_h,b_h]$，并在 $I$ 上令 $w(t)=(t-a_h)^2(b_h-t)^2$，其余处置零。
在 $I$ 上取 $v_h=-w$，按 $v_h(1-t)=-v_h(t)$ 延拓，在两个支集之外置零。
这把暴露的正臂缩短，同时把互补臂延长相同的长度。
新互补臂仍在 $B_a(1-t)$ 以下，新增的整个三角形都被共同尾部遮住。
原递增邻域中的一个紧高半径带严格损失端点相位，正如式~\eqref{eq:upper-height-cost} 中的论证。
选中分支的导数在紧内部区间上严格为负，所以充分小的 $C^1$ 改动保留递减性和起始条件。
第一种情形由此完成，全局端点映射递增并非所需假设。

% [U036]
若 $r=p_0$，则 $a-s$ 及其一阶导数都在 $\alpha$ 和 $1-\alpha$ 处为零。
旧递增映射 $g_{p_0,s}$ 因而为 $g_{p_0,a}$ 提供两个局部递增邻域。
在每个邻域内选一个更早的截断点 $b$；对于所有更早的源，左侧有
\[
 g_{p_0,a}(u)\le g_{p_0,s}(u)\le g_{p_0,s}(b)
                          <g_{p_0,s}(\alpha)
\]
右侧在 $1-\alpha$ 处有相应的不等式。
将两个源支集选得足够靠近各自边界，使输出窗口既避开所有这些早期源的相位界，又彼此分离。

% [U037]
对支撑在这两个区间内的任意紧支集反射 $C^1$ 凸起 $v_h$，辅助放置 $q_\eta=p_0+\eta v_h$ 在小 $\eta$ 下仍保持这些严格相位间隙；检查后一个窗口时，较早的已修改窗口也受同一控制。
对未改变的源，若 $r\le p_0$，可与这个辅助映射比较；若 $r>p_0$，中位数给出 $r\le B_s\le B_a$，所以整条臂都被遮住。
更晚的合法源严格低于首次到达半径，因为 $a(u)/(u-z)$ 在首次到达后递减。
这些情形穷尽了更早、更晚及被遮住的源，证明新旧集合在两个窗口上的实际径向上边界均为 $\max(R,T_a)$。
这个辅助放置只承担相位比较，不要求 $(p_0,a)$ 在全局满足起始条件。

% [U038]
接触刚性仍在局部适用：每个高度分片的导数次数是十五，而 $p_0$ 的次数是十九。
实现阈值最大值的源参数随 $t$ 严格递增，因此可先避开其有限多个分片接缝，再使用定理~\ref{thm:upper-polynomial} 的有理复合证明。
令 $D_h=\partial_x\Psi^{a}_t(p_0(t))$，其中 $\Psi^{a}$ 是以高度 $a$ 构成的配对铰链。
选取避开接触点和 $D_h$ 零点的紧激活区域 $I=[a_h,b_h]$。
在 $I$ 上令 $w(t)=(t-a_h)^2(b_h-t)^2$，区间外置零；在该区间取 $v_h=-\operatorname{sign}(D_h)w$，再像第一种情形那样反射并补零。
每条被改变的完整底边条带，其相位介于新旧端点相位之间；随后径向填充不改变相位。
因此，局限在两个已验证窗口内的是全部实际几何变化，而不仅是端点变化。

% [U039]
在激活窗口上有 $p^2g'=p^2-pa'+ap'$，分部积分消去了凸起的导数。
把两个窗口配对，得到精确有限差分
\[
 C_{\mathrm{phys}}(r+\eta v_h,a)-C_{\mathrm{phys}}(r,a)
  =\pi(\eta L_h+\eta^2Q_h),\quad
 L_h=-\int_I|D_h|w<0,\quad Q_h=\int_I\nu w^2>0,
\]
其中按激活情形，$D_h=2(p_0-a')$ 或 $4p_0-2-4a'$，相应地 $\nu=1$ 或 $2$。
选一个正的 $\eta_h<-L_h/Q_h$，同时满足此前固定的起始、导数、相位间隙与激活区域裕量。
第一种情形改用紧尾部与起始邻域的裕量选 $\eta_h$，其正节省已由严格半径带论证保证。
两种情形都先固定 $p_*$ 和 $c_{\mathrm{hl}}$，然后再应用命题~\ref{prop:upper-recovery}。
\end{proof}

\subsection{在原始微观尺度上归一化整条弦}
\label{subsec:upper-chords}

% [U040]
一次性固定命题~\ref{prop:upper-longitudinal} 中的轮廓对 $(p_*,a)$，令 $g(t)=t-a(t)/p_*(t)$，并将其原始有限恢复简记为 $K_{n,k}$。
接下来的操作同时使用这个集合、它的原始记录以及 $\delta_n=\pi/n$。
这个操作依赖原始记录，仅凭未标记的集合不能确定。

% [U041]
对每条原始闭记录 $e=(i,q)$，记
\[
 N_e=[l_e,r_e]u_q+h_e n_q,\qquad r_e-l_e=1.
\]
在所有长度 $\ell=U-L\ge1$、整段包含于集合内的区间 $C=[L,U]u_q+Hn_q\subset K_{n,k}$ 中，施加约束
\begin{equation}\label{eq:upper-Q}
 (L-l_e)^2\le\ell-1,\qquad (U-r_e)^2\le\ell-1,
 \qquad \left(\frac{H-h_e}{\delta_n}\right)^2\le\ell-1.
\end{equation}
将 $\ell$ 最大化，对每条记录保留\emph{所有}极大解，并令
\begin{equation}\label{eq:upper-Q-output}
 S^Q_{n,k}=\overline{\bigcup_e\ \bigcup_{C\ \text{为记录 }e\text{ 的极大解}}
                              \operatorname{conv}(0,C/\ell)}.
\end{equation}
包含约束逐点检查整条区间；不连通直线截面中的两段，不能跨过孔洞合成一条弦。
规则只对旧集合执行一次，不按有利源筛选，也不迭代。

\begin{theorem}[归一化约束延长给出的固定实际面积删除]
\label{thm:upper-Q}
存在固定 $\Gamma_{\mathrm{ms}}>0$ 及整数 $N$，使每个奇数 $n\ge N$ 都有自己的 $k_0(n)$，并且所有原始细分 $k\ge k_0(n)$ 均满足
\[
 S^Q_{n,k}\subset K_{n,k},\qquad
 |K_{n,k}\setminus S^Q_{n,k}|\ge\Gamma_{\mathrm{ms}}.
\]
特别地，
\begin{equation}\label{eq:upper-Q-limit}
 \limsup_{\substack{n\to\infty\\n\ \mathrm{odd}}}
 \limsup_{k\to\infty}|S^Q_{n,k}|
                  \le c_{\mathrm{hl}}-\Gamma_{\mathrm{ms}},
 \qquad A_*\le c_{\mathrm{hl}}-\Gamma_{\mathrm{ms}}.
\end{equation}
$\Gamma_{\mathrm{ms}}$ 的一个精确定义是证明中所给正带积分的一半。
\end{theorem}

\begin{proof}
% [U042]
旧线段本身可行，紧性又使所含区间的坐标和长度全部有界。
整段包含关系在端点收敛下是闭的，式~\eqref{eq:upper-Q} 的三个约束也都是闭的，因此最大值存在。
每个 $C/\ell$ 的长度恰为一，并由星形性知它包含于 $K_{n,k}$；它的整个三角形以及这些三角形并集的闭包也都包含于其中。
方向、零高度记录、射影重提升以及并列极大解，一个也没有删去。
若 $\ell=1$，约束就迫使 $C=N_e$ 精确成立。

\paragraph{在每个实现上边界的源附近延长整条弦。}
% [U043]
在已选纵向激活区域内取嵌套紧区间，使较小区间包含某个紧的正则暴露相位区间 $Z\subset(0,1)$ 上的全部首次到达源参数。
旧支集端点满足 $r>d$，$p_*$ 在那里未变且严格递减，所以整个支集上都有 $p_*>d$。
其局部径向上边界 $R(z)=p_*(g^{-1}(z))$ 因而高于 $d$，并与共同尾部及非等价的相位表示分离。
此处只对已证明的局部首次到达映射取逆，并未假设全局可逆。
我们将在该上边界的\emph{每个}等价源表示的邻域上得到固定的长度下界。

% [U044]
互补臂被遮住时，它的端点严格位于实际尾部上边界之下。
单激活的多项式情形也包括在内。
尾部的极大源起始可容许，严格的尾部支配同时给出底边内部裕量和径向收缩裕量。
这给出一段真正连通的向外延续；它及其与旧线段的连接处，每一点都严格位于某个源三角形内部。
保留原来的整条线段，在同一支撑直线上接上一段足够短的这种延续段。
长度增益随延长量线性变化，约束中的纵向误差平方随延长量二次变化，法向误差为零，所以式~\eqref{eq:upper-Q} 带有严格裕量地成立。

% [U045]
双激活时，旧的两个端点都暴露在外，所需的较长区间必须移到邻近的支撑直线上。
在偶数单元内，源映射为
\[
 F_n(t,x)=\operatorname{Rot}_{\delta_n t}(x,-\delta_n a(t)),
                         \qquad p_*(t)-1\le x\le p_*(t).
\]
在 $t$ 处的坐标系中，相邻源 $u$ 的正端点坐标为
\begin{align}
 X_+(u;t)&=p_*(u)\cos(\delta_n(u-t))
                    +\delta_n a(u)\sin(\delta_n(u-t)),\\
 A_+(u;t)&=a(u)\cos(\delta_n(u-t))
                    -p_*(u)\frac{\sin(\delta_n(u-t))}{\delta_n},
\end{align}
其法向坐标为 $-\delta_n A_+$。
负端点的公式把 $p_*$ 换成 $p_*-1$。
在 $u=t$ 处，它们的导数对为
\[
 (X_{+,u},A_{+,u})=(p_*'+\delta_n^2a,a'-p_*),\qquad
 (X_{-,u},A_{-,u})=(p_*'+\delta_n^2a,a'+1-p_*).
\]
在选定的紧激活区域上，$p_*'<0$、$p_*>a'$、$a'+1-p_*>0$ 都具有严格裕量。
所以，对小的正 $e$，将支撑线移至 $-\delta_n(a(t)+e)$ 会使正臂一侧的交点右移、负臂一侧的交点左移。
两端已经延长。内部是否有孔洞，还须单独检查。

% [U046]
精确的源参数--底边坐标 Jacobian 的绝对值为 $\delta_n|x-a'(t)|$。
避开唯一折叠点 $x=a'(t)$，源映射的像是开的；把目标点略向径向外侧移动，用相邻底边内点表示，再沿同一射线缩回目标点，就使目标点严格落入一个源三角形内部。
这个直角坐标论证也处理 $x=0$，无需除以纵向坐标。
折叠点本身还需要另一个见证。

% [U047]
在左侧支集的折叠点处，令 $\rho=a'(t)>0$、$z=a(t)/\rho-t=L_a(\rho)$。
严格凹性给出 $T_a(z)=\rho$。
高度构造使 $\rho$ 落在受保护窗口内，并给出 $G_{r,a}(\rho)>L_a(\rho)$，因此折叠点严格位于高度阶段集合的内部。
要把这个事实传到实际的 $p_*$ 集合，使用命题~\ref{prop:upper-longitudinal} 已证明的完整条带定位：在两个改动窗口外，集合不变；在双激活窗口内，新上边界仍严格高于 $T_a$；在窗口边界，新旧数据相同。
这处理了折叠相位及其带符号等价表示，而不要求全局成立 $G_{p_*,a}>L_a$。

% [U048]
还需逐点证明：每个折叠点严格落在某一个源三角形内。
在稍大的半径和仍低于端点上边界的邻近相位处，激活前缀中有一个源，其端点相位严格超过折叠相位。
从 $g(0)=0$ 出发，连续性给出一次交叉，再给出附近的一个源，其底边坐标严格介于折叠半径与该源端点半径之间。
这个底边内点位于折叠射线上，严格径向收缩后到达折叠点。
因此，每个折叠点都属于一个实际正高度源三角形的内部，而不仅属于并集的内部。
由紧性可取有限多个这样的三角形，覆盖整个闭折叠邻域。

% [U049]
折叠邻域、非折叠中部以及两个正则端点坐标片的内侧，覆盖了整条平移后的直线区间。
略微裁去每个新端点，同时保留两端正的延长量。
遍历固定紧邻域内的所有源、所有充分大的奇数 $n$ 及 $0<e\le e_0$，定义 $c_\ell$ 为 $(\ell(e)-1)/e$ 的下确界，定义 $C_Q$ 为一与下列三个量的上确界之最大值：
\[
 |L(e)-l|/e,\qquad |U(e)-r|/e,\qquad |H(e)-h|/(\delta_n e).
\]
选取足够小的 $e_0$ 后，端点导数的严格界给出 $c_\ell>0$，一致隐函数估计给出 $C_Q<\infty$。
选一个满足 $C_Q^2e<c_\ell$ 的 $e\in(0,e_0)$，令 $\epsilon_\ell=c_\ell e/2>0$。
必要时降低共同下界，以同时纳入接续尾部的情形；这样每个相关记录邻域都有一个合格的完整区间，长度至少为 $1+\epsilon_\ell$。
位移和长度常数只用于证明可行长区间存在，不进入极大化规则的参数。

% [U050]
所有严格裕量都在法向坐标除以 $\delta_n$ 后的坐标系中成立。
对应的有限子覆盖对所有充分大的 $n$ 持续存在；随后对每个固定 $n$，在所有充分细的\emph{原始} $(b,A)$ 同步冻结下也持续存在。
两侧闭标签都满足同一组一致端点误差和归一化高度误差界。
接续尾部的情形按定义保留原线段，只有连接处和新增部分需要严格见证。
平移情形则为裁短后的整个区间提供这种见证。
有限覆盖已经足够；不必假设最大弦长连续，也不必假设直线连通分支下半连续。

\paragraph{归一化约束阻止重新填回。}
% [U051]
先考虑旧三角形中的一个点列，其极限半径为正，归一化相位位于 $Z$。
把每个点写成 $\lambda(xu_q+\delta_n A_i n_q)$。
共同外界和正半径使 $|x|$ 与零保持距离，$A_i$ 又有界，因此相对的带符号相位等价表示只有有限种可能。
完整底边数据的一个子列收敛到合法的极限底边及径向参数。
因此，旧集合点列的极限仍在完整极限集合内；这里需要的比面积收敛更多。
若极限点达到 $R(z)$，尾部分离、早期源相位间隙以及所有更晚合法半径的严格递减，共同迫使它成为 $\lambda=1$ 的首次到达端点。
这个端点无论由哪个等价源表示，都落在一个已经证明具有长度下界的邻域内。

% [U052]
假设沿奇数 $n_j\to\infty$，每个 $n_j$ 处的原始网格都能选得任意细，却仍有重新生成的点趋近旧上边界。
对新三角形中的点 $y$，与之关联的极大长度满足 $\ell y\in K_{n,k}$。
因此，旧集合的极限包含关系与径向极大性迫使 $\ell\to1$。
式~\eqref{eq:upper-Q} 随之迫使
\[
 L-l_e\to0,\qquad U-r_e\to0,\qquad
                  (H-h_e)/\delta_n\to0.
\]
最后一个极限比实际位移趋于零更强，它使新旧归一化底边数据一致。
极限数据实现旧上边界，因此相应原始记录最终落入极大长度至少为 $1+\epsilon_\ell$ 的邻域，矛盾。
对式~\eqref{eq:upper-Q-output} 中的闭包点，先用实际新三角形点以 $o(\delta_n)$ 的误差逼近它，再用同一论证。
这个矛盾不依赖并列极大解的选择，并按定理中的量词顺序给出固定的归一化排除邻域。

\paragraph{缺失带具有不消失的实际面积。}
% [U053]
选正宽度闭区间 $Z'\subset\operatorname{int}Z$ 及 $0<\sigma_1<\sigma_2$，使紧带
\[
 \mathcal B=\{(z,\rho):z\in Z',\ R(z)-\sigma_2\le\rho\le R(z)-\sigma_1\}
\]
位于上述排除邻域内，且所有半径都大于 $d$。
带内每个点都有严格的旧三角形见证：将首次到达源略向后移，在同一射线上取此时已经位于底边内部的坐标，再径向收缩。
有限子覆盖保证它的所有重复有限角度副本都在原始 $K_{n,k}$ 中持续存在。
这些副本与整个闭集 $S^Q_{n,k}$ 不交。
先合并标签及等价相位表示，再按原来的双侧复制密度，只对一个实际相位带计面积，定义
\begin{equation}\label{eq:upper-gamma-definition}
 D_{\mathrm{ms}}=\pi\int_{Z'}\int_{R(z)-\sigma_2}^{R(z)-\sigma_1}
                    \rho\,d\rho\,dz>0,
 \qquad \Gamma_{\mathrm{ms}}=D_{\mathrm{ms}}/2.
\end{equation}
略去有界个数的全局首尾单元，所改变的面积趋于零。
因此，先取所有充分大的奇数 $n$，再取所有充分细的原始网格，这些副本的实际面积均超过 $\Gamma_{\mathrm{ms}}$。
这证明了面积删除及嵌套上极限不等式。

% [U054]
还须把输出恢复为原始有限竞争体。每个 $S^Q_{n,k}$ 都紧、以原点为星心，并在每个方向含完整单位线段，故引理~\ref{lem:finite-open-recovery} 给出面积至多为 $|S^Q_{n,k}|+\varepsilon$ 的原始有限体。
这里先固定奇数 $n$，再固定它的原始细分 $k$，最后令独立的恢复误差 $\varepsilon$ 趋于零；随后才使用已经证明的先内层 $k$、后外层 $n$ 的嵌套上极限。
这证明式~\eqref{eq:upper-Q-limit} 的最后一个不等式，无需假设输出面积极限存在，也无需任意对角收敛定理。
\end{proof}

\subsection{连续针运动给出的独立改进}
\label{subsec:upper-coherent}

% [U055]
如果一个原始有限体包含下列完整线段族，就称它\emph{相干}：
\[
 c(q)+[-1/2,1/2]u_q\qquad(q\in\mathbb {RP}^1)
\]
这里 $c:\mathbb {RP}^1\to\mathbb R^2$ 是某个连续的实际中点映射。
$A_{\mathrm{coh}}$ 是这个原始有限相干子类上的面积下确界。
保留全部极大解并未给出连续的极大解选择；要得到连续运动，必须另构造一个输出族。

\begin{proposition}[相干的正面积节省]\label{prop:upper-coherent}
对于同一个固定连续源 $(p_*,a)$，存在以充分大奇数 $n$ 为指标的原始有限相干体 $G_n$，以及固定 $\Delta_{\mathrm{coh}}>0$，使得
\[
 \limsup_{\substack{n\to\infty\\n\ \mathrm{odd}}}|G_n|
       \le c_{\mathrm{hl}}-\Delta_{\mathrm{coh}},\qquad
 A_*\le A_{\mathrm{coh}}\le c_{\mathrm{hl}}-\Delta_{\mathrm{coh}}<c_{\mathrm{hl}}.
\]
这个构造既不比较 $\Delta_{\mathrm{coh}}$ 与 $\Gamma_{\mathrm{ms}}$，也不把两种节省相加。
\end{proposition}

\begin{proof}
% [U056]
在每个实现旧上边界的源附近，把前述较长区间的构造写成从原线段出发的连续路径 $C_n(t,e)$。参数及长度记为
\[
 0\le e\le e_0,\qquad \ell_n(t,e)=\operatorname{length}C_n(t,e).
\]
互补臂被遮住时，沿已经证明连通的尾部恰好接续长度 $e$。
双激活时，解精确的有限角度端点方程 $A_\pm(u;t)=a(t)+e$，将所得右、左交点坐标记为 $R_{\mathrm{int}}$、$L_{\mathrm{int}}$。
采用支撑线 $-\delta_n(a(t)+e)$，并取保留一半延长量的端点
\[
 U_n(t,e)=p_*(t)+\frac{R_{\mathrm{int}}(t,e)-p_*(t)}2,\qquad
 L_n(t,e)=p_*(t)-1+
                  \frac{L_{\mathrm{int}}(t,e)-(p_*(t)-1)}2.
\]
一致隐函数估计使这些路径对所有小 $e$ 都存在，也包括 $e=0$。
弦证明中固定的折叠邻域与中部邻域覆盖所有中间位移；精确端点坐标片则保证，对每个 $e>0$，保留一半延长量的两端仍在内部。
在 $e=0$ 处，直接使用原线段。
这样就不必把有限角度近似误差除以一个可能任意快地趋于零的位移。

% [U057]
选连续截断函数 $\chi$，在每个实现极限上边界的源附近等于一，并在到达路径支集边界和任何整数接缝之前降为零。
通过等价的带符号源表示传递同一个截断函数。
在这些邻域内使用 $C_n(t,e_0\chi(t))$，其余处使用原线段，再以原点为中心，将每个选定区间按其自身长度缩放。
所得单位族 $V_n(q)$ 的实际中点连续，并生成紧的单位生成星形集 $S_n\subset E_n[p_*,a]$。
在射影接缝处，奇数 $n$ 使 $b$ 变号，坐标系也同时变号，所以实际中点闭合；截断函数不改变这个接缝。
每个选定区间都满足式~\eqref{eq:upper-Q}。
在固定的紧源邻域内取遍所有源，并取遍所有充分大奇数 $n$，定义 $\epsilon_c$ 为 $\ell_n(t,e_0)-1$ 下确界的一半；一致路径估计保证它为正。
每个边界实现源附近，\emph{实际选中}的长度都超过 $1+\epsilon_c$；仅有可选最大长度达到这个下界还不够。

% [U058]
不作极大化，同一反证仍成立：假想回填迫使缩放长度趋于一，归一化约束使极限数据实现旧上边界，而实际选中长度的下界与这个极限矛盾。
按式~\eqref{eq:upper-gamma-definition} 选取一个被删除的紧内部带 $\mathcal B_c$，以 $D_c>0$ 记其相应实际积分。
固定恢复前的删除预算 $\Delta_0=D_c/2$，并定义 $\Delta_{\mathrm{coh}}=D_c/4$。
对所有充分大的奇数 $n$，$E_n\setminus S_n$ 的实际面积至少为 $\Delta_0$。
这个删除带来自新的连续构造，不能与定义 $D_{\mathrm{ms}}$ 的带视为同一个带。

% [U059]
要让有限恢复保留这条运动，恢复体必须包含原集合；面积相近本身并不够。
令 $\mu=\min_{[0,1]}\min(p_*,1-p_*)>0$、$L_*=\max_{[0,1]}\max(p_*,1-p_*)$，并令 $R_*=(1+L_*)/2<1$。
界 $\sqrt{L_*^2+\delta_n^2\|a\|_\infty^2}$ 说明，所有充分靠后的连续集合都包含在半径 $R_*$ 的圆盘中。
每个选中的长区间都延长原纵向端点，且长度至多为 $2R_*$，所以归一化后的两臂都至少为 $\beta_c=\mu/(2R_*)>0$。
若实际中点为 $v_n(q)$，则 $[0,2\pi]$ 上的端点 $v_n(q)+u_q/2$ 的辐角为
\[
 q+\arctan\frac{v_n(q)\cdot n_q}{1/2+v_n(q)\cdot u_q}.
\]
分母为正，修正项连续且以 $2\pi$ 为周期。
所以辐角映射的度为一，并在至少为 $\beta_c$ 的半径处到达每条实际射线。
星形性给出内含关系 $B_{\beta_c}\subset S_n$，并未把这个圆盘中的每一点都证明为星心。

% [U060]
设紧的原点星形集 $S$ 包含 $B_\beta$，在每个方向都含完整单位线段，并且每个非零点都位于某条这样的单位线段与原点生成的三角形内。则它允许有限包含恢复，且面积误差与外邻域误差都可任意小。
为此，先按 $\lambda\in(1,2)$ 放大，再把每条放大后的生成线段分成靠两端的两个相互重叠的单位子线段。
它们的并集是整条放大线段，各自的窄角旋转扇形则是位于 $\lambda S+B_\eta$ 内的原始常数记录扫掠。
$S$ 的每个非零点都严格位于某个放大底边点的径向内侧，所以这些扇形为它提供开空间邻域。
对 $S\cap\{|x|\ge\beta/2\}$ 取有限子覆盖，先覆盖这个环形部分，再以径向填充覆盖 $B_{\beta/2}$。
再另加一个有限方向覆盖，确保记录区间覆盖整个 $\mathbb {RP}^1$。
当 $\lambda\downarrow1$、$\eta\downarrow0$ 时，共同紧外包集合的面积趋于 $|S|$，一次性控制所有扇形。
因此，对任意指定的 $\tau,\varepsilon>0$，可得到 $S\subset G\subset S+B_\tau$ 及 $|G|\le|S|+\varepsilon$。

% [U061]
对每个固定 $S_n$ 应用这个构造，并使面积误差趋于零。
原始有限体 $G_n$ 包含原来的整个连续族 $V_n$，相干性因而保留。
恢复前的删除量给出所列上界，其中使用上述较小的 $\Delta_{\mathrm{coh}}$。
这里没有从 $S^Q_{n,k}$ 内部推断出连续选择函数。
\end{proof}

\subsection{网格轴任意短的有限竞争体}
\label{subsec:upper-short-axes}

% [U062]
改进后的源也为连续运动的几何障碍提供具体的有限输入。
这里沿用已经证明面积极限为 $c_{\mathrm{height}}<U$ 的高度阶段序列，不给归一化弦输出附加其他性质。
下面的命题给出短轴障碍定理所需的固定集合。

\begin{proposition}[一个原始有限竞争短轴体]
\label{prop:upper-short-axis-example}
存在原始有限体 $K$、半径 $R_0<1$、整数 $m\ge3$ 以及数
\[
 0<a_0<\min(1-R_0,1/2)
\]
使 $|K|<U$、$K\subset B_{R_0}$，且对每个 $\ell=0,\ldots,m-1$ 都有
\[
 K\cap\mathbb R u_{\ell\pi/m}\subset B_{a_0}.
\]
下面的有限窗口替换和指标选择将指定一个满足条件的集合。
\end{proposition}

\begin{proof}
% [U063]
令 $n_j=2j+1$。取最小细分指标 $k_j$，要求从它开始的每个后续细分，在固定 $n_j$ 下的面积误差都小于 $1/j$。
命题~\ref{prop:upper-recovery} 保证这个指标有限，并给出 $K_j=K_{n_j,k_j}[r,a]$、$\delta_j=\pi/n_j\downarrow0$ 及 $|K_j|\to c_{\mathrm{height}}$。
定义
\[
 L=\max_{[0,1]}\max(r,1-r)<1,\qquad M=\max_{[0,1]}a,
 \qquad R_0=(1+L)/2<1.
\]
每个原始底边坐标的纵向绝对值至多为 $L$，高度至多为 $M\delta_j$。
所以其范数至多为 $\sqrt{L^2+M^2\delta_j^2}$，每个充分靠后的 $K_j$ 都位于同一个 $B_{R_0}$ 中。
这个界同时控制实际冻结记录及接缝两侧的副本。

% [U064]
令 $\gamma_j=\sqrt{\delta_j}$、$m_j=\lfloor\delta_j^{-1/4}\rfloor$。
在每个网格方向 $q_\ell=\ell\pi/m_j$ 处，只从旧记录中删去射影半径为 $\gamma_j$ 的开方向窗口。
保留旧记录与闭补集的交，再加入整个闭替换窗口
\[
 [-1/2,1/2]u_{q_\ell+t}+\tan(\gamma_j)n_{q_\ell+t},
                         \qquad -\gamma_j\le t\le\gamma_j.
\]
以 $F_j$ 表示所有保留三角形与替换三角形的并集，它们均以原点为顶点。
对大的 $j$，有 $m_j\ge3$、$\gamma_j<\pi/(4m_j)$，所以这些小方向窗口互不相交，所有闭端点也都保留。
在射影接缝处分割时同时反转符号，便严格得到原始有限记录形式，方向覆盖也是精确的。
替换底边的范数至多为 $1/2+\tan\gamma_j$，所以对所有充分大的 $j$，同一个 $F_j$ 也位于 $B_{R_0}$ 中。

% [U065]
来自保留旧三角形、落在网格轴上的点，其半径至多为 $M\delta_j/\sin\gamma_j$。
事实上，每个保留源与该轴的方向距离至少为 $\gamma_j$，其法向坐标又受 $M\delta_j$ 控制。
该轴自身的替换片只在原点与轴相交，因为每个底边点的法向坐标至少为
\[
 -\tfrac12\sin\gamma_j+\tan\gamma_j\cos\gamma_j
                                      =\tfrac12\sin\gamma_j>0.
\]
对其他每个替换片，源方向与目标轴的射影距离至少为 $\pi/m_j-\gamma_j$。将高度 $\tan\gamma_j$ 除以这个距离下界的正弦，就得到相应的径向上界。
因此，\emph{整个}修补后的集合都满足轴界，其中
\begin{equation}\label{eq:upper-short-axis-radius}
 a_j=\max\left\{\frac{M\delta_j}{\sin\gamma_j},
              \frac{\tan\gamma_j}{\sin(\pi/m_j-\gamma_j)}\right\}\longrightarrow0.
\end{equation}
所有其他替换片可能造成的回填，以及每条网格直线的正负两条射线，都已计入这个界。

% [U066]
一个替换片实际新增的部分位于面积至多为下式的矩形中：
\[
 E(\gamma)=(1+2\tan\gamma)
                           (\tfrac12\sin\gamma+\tan\gamma).
\]
一个窗口中被移除的旧部分，相对于中央轴的法向位移至多为 $R_0\sin\gamma_j+M\delta_j$，因此对应矩形面积至多为 $4R_0(R_0\sin\gamma_j+M\delta_j)$。
于是
\[
 |F_j\triangle K_j|\le m_jE(\gamma_j)
       +4m_jR_0(R_0\sin\gamma_j+M\delta_j)\longrightarrow0,
\]
这里使用 $m_j\gamma_j\to0$ 及 $m_j\delta_j\to0$。
特别地，$|F_j|\to c_{\mathrm{height}}<U$。
越过保留的初始段后，取 $j_0$ 为首个使下列条件同时成立的指标：
\[
 F_j\subset B_{R_0},\quad |F_j|<U,\quad m_j\ge3,
 \quad 0<a_j<\min(1-R_0,1/2)
\]
并令 $K=F_{j_0}$、$m=m_{j_0}$、$a_0=a_{j_0}$。
最终成立性保证这是一项确定的有限构造；证明并未算出记录数，也没有给出有效指标。
\end{proof}

% [U067]
同一个替换操作还给出更强的结构性序列：$F_j$ 的所有凸子集的面积一致趋于零。
否则，在共同圆盘内可取面积有正下界的紧凸子集序列，其 Hausdorff 极限仍是正面积凸集。
在极限内部取一个较小圆盘，再用有限多边形包围它；逐顶点逼近多边形并用凸性填满，就知这个圆盘仍在充分靠后的凸近似内。
足够密的网格射线会在与零保持正距离的半径处穿过这个圆盘，与式~\eqref{eq:upper-short-axis-radius} 矛盾。
因此，即使原始有限竞争体的面积低于 $U$，较大的凸块也未必保留。
这个结论使用凸集内部的持续性，对一般非凸 Hausdorff 极限不作同样断言。

\subsection{精确固定数据与尚未建立的变分联系}
\label{subsec:upper-evidence}

% [U068]
冻结输入由下面两个有序元组精确给定，无须小数拟合。
若 $(\alpha_i)_{i=0}^{19}$ 和 $(s_i)_{i=0}^{16}$ 按所列顺序取值，则
$\alpha_{\mathrm{fr}}(t)=\sum_{i=0}^{19}\alpha_i t^i$，$s(t)=\sum_{i=0}^{16}s_i t^i$。
\begingroup\small
\begin{align*}
(\alpha_i)={}&\bigl(99/280,-840817/1260000,-1065631/900000,636709/50000,\\
 &-1046431/28125,-41037553/56250,88036453/6250,-441121411/3125,\\
 &2952524356/3125,-13982764696/3125,48196297776/3125,\\
 &-123074065696/3125,235220271104/3125,-337443094528/3125,\\
 &361517023232/3125,-284890759168/3125,160289521664/3125,\\
 &-60934848512/3125,14025752576/3125,-1476395008/3125\bigr),\\[2mm]
(s_i)={}&\bigl(0,107161751/128520000,-102943399/642600000,\\
 &-1107847/105000,1266525811/9450000,-899826097/590625,\\
 &54998333/4375,-41734226272/590625,162045709624/590625,\\
 &-2352117376/3125,4626917056/3125,-6552918144/3125,\\
 &6626537152/3125,-4668260352/3125,2176843776/3125,\\
 &-603979776/3125,75497472/3125\bigr).
\end{align*}
\endgroup

% [U069]
这些元组固定了本节的全部多项式输入。附录~\ref{sec:formal-evidence} 将所用证据分为系数与恒等式、严格符号与起始条件、端点映射、冻结放置单调性、固定带区间支配、端点与尾部界，以及历史面积比较。随附的《上界证据映射》逐项列出数学陈述和精确源码声明。\footnote{电子附件：\texttt{\detokenize{provenance/UPPER_EVIDENCE_MAP.md}}。}
其中，$p_0$ 的非递增性使用冻结放置导数非正的定理，不以导数绝对值界代替；首段、末段支配与半相位包围也各有自己的认证范围。
历史基准证书另行提供 $c_{\mathrm{fr}}<U$；系数文件中单独命名的目标面积，既不是这里的积分，也不能代替 $U$。

% [U070]
以上所有新区间和振幅，都从由紧集上严格不等式定义的非空集合中选出。
固定任意一组这样的选择，就确定了完整实际集合族及其代价；它们都不是新求值的系数表或数值极小解。
这些选择不仅存在，而且在每个复杂度极限之前，就已固定了正的面积节省。
这已经足以给出严格上界比较，而节省的数值大小仍未经求值。

% [U071]
铰链比较仍以严格凹性、反射、起始条件及交替放置类为假设。
尚未证明能把每个原始有限体都送入这个类、且不增加面积的映射。
即使在类内，参考补全选择函数能否取等号，仍归结为前述激活折叠问题。
用冻结轮廓在每个半径上的径向占据作为逐半径下界，这个设想是错误的：合法圆盘 $B_{1/2}$ 在 $\rho>1/2$ 时没有截面，而冻结正臂在这类半径的一个开区间上给出非常值相位区间。
积分比较允许不同半径的占据相互补偿，因而不同于已被否定的逐半径比较；这样的积分断言仍未建立。

% [U072]
有限可达性问题也仍然开放：现有结果没有证明每个原始有限体都能严格改进。
这些下降针对指定的源或保留记录信息的输出；短轴构造给出的是一个特定竞争体，而非普适标准形。
同样，相干正节省证明的是 $A_{\mathrm{coh}}$ 的上界，并未证明 $A_{\mathrm{coh}}=A_*$，也未证明两个下确界严格分离。
某个固定短轴体的相干逼近障碍可以迫使添加面积，但另一个相干竞争体可能删去旧部分，抵消添加的面积。
全类问题还缺两端在同一数值上的匹配：一端是普适的积分下界比较，另一端是实际实现该值的序列。
